% Options for packages loaded elsewhere
\PassOptionsToPackage{unicode}{hyperref}
\PassOptionsToPackage{hyphens}{url}
\documentclass[
  ignorenonframetext,
]{beamer}
\newif\ifbibliography
\usepackage{pgfpages}
\setbeamertemplate{caption}[numbered]
\setbeamertemplate{caption label separator}{: }
\setbeamercolor{caption name}{fg=normal text.fg}
\beamertemplatenavigationsymbolsempty
% remove section numbering
\setbeamertemplate{part page}{
  \centering
  \begin{beamercolorbox}[sep=16pt,center]{part title}
    \usebeamerfont{part title}\insertpart\par
  \end{beamercolorbox}
}
\setbeamertemplate{section page}{
  \centering
  \begin{beamercolorbox}[sep=12pt,center]{section title}
    \usebeamerfont{section title}\insertsection\par
  \end{beamercolorbox}
}
\setbeamertemplate{subsection page}{
  \centering
  \begin{beamercolorbox}[sep=8pt,center]{subsection title}
    \usebeamerfont{subsection title}\insertsubsection\par
  \end{beamercolorbox}
}
% Prevent slide breaks in the middle of a paragraph
\widowpenalties 1 10000
\raggedbottom
\AtBeginPart{
  \frame{\partpage}
}
\AtBeginSection{
  \ifbibliography
  \else
    \frame{\sectionpage}
  \fi
}
\AtBeginSubsection{
  \frame{\subsectionpage}
}
\usepackage{iftex}
\ifPDFTeX
  \usepackage[T1]{fontenc}
  \usepackage[utf8]{inputenc}
  \usepackage{textcomp} % provide euro and other symbols
\else % if luatex or xetex
  \usepackage{unicode-math} % this also loads fontspec
  \defaultfontfeatures{Scale=MatchLowercase}
  \defaultfontfeatures[\rmfamily]{Ligatures=TeX,Scale=1}
\fi
\usepackage{lmodern}
\ifPDFTeX\else
  % xetex/luatex font selection
\fi
% Use upquote if available, for straight quotes in verbatim environments
\IfFileExists{upquote.sty}{\usepackage{upquote}}{}
\IfFileExists{microtype.sty}{% use microtype if available
  \usepackage[]{microtype}
  \UseMicrotypeSet[protrusion]{basicmath} % disable protrusion for tt fonts
}{}
\makeatletter
\@ifundefined{KOMAClassName}{% if non-KOMA class
  \IfFileExists{parskip.sty}{%
    \usepackage{parskip}
  }{% else
    \setlength{\parindent}{0pt}
    \setlength{\parskip}{6pt plus 2pt minus 1pt}}
}{% if KOMA class
  \KOMAoptions{parskip=half}}
\makeatother
\setlength{\emergencystretch}{3em} % prevent overfull lines
\providecommand{\tightlist}{%
  \setlength{\itemsep}{0pt}\setlength{\parskip}{0pt}}
\usepackage{bookmark}
\IfFileExists{xurl.sty}{\usepackage{xurl}}{} % add URL line breaks if available
\urlstyle{same}
\hypersetup{
  hidelinks,
  pdfcreator={LaTeX via pandoc}}

\author{\texorpdfstring{}{}}
\date{}

\begin{document}

\section{Measure-theoretic context
(qualitative)}\label{measure-theoretic-context-qualitative}

\begin{frame}{Almost every number admits sharp approximations}
\protect\phantomsection\label{almost-every-number-admits-sharp-approximations}
Dirichlet's theorem (pigeonhole principle on fractional parts; see
{[}@cassels1957{]}) implies that for every irrational \(x\) and every
\(Q>1\) there exist integers \(p,q\) with \(1 \leq q \leq Q\) and
\(|qx-p|<1/Q\), hence {[} \left\textbar x-\frac{p}{q}\right\textbar{}
\textless{} \frac{1}{qQ} \leq \frac{1}{Q}. {]} So
\(\delta_Q(x) \leq 1/Q\) always for irrational \(x\); the
\textbf{leading} scale is universal. The interesting structure at
moderate \(Q\) lies in \textbf{how much smaller} than \(1/Q\) the
minimum can be---left-tail behaviour of \(\delta_Q(X)\) under a
reference law---and in the \textbf{relative} placement of named
constants inside that tail.
\end{frame}

\begin{frame}{Badly approximable numbers}
\protect\phantomsection\label{badly-approximable-numbers}
An irrational \(x\) is \emph{badly approximable} if there exists \(c>0\)
such that \(|x-p/q|>c/q^2\) for all rationals \(p/q\). Equivalently, the
partial quotients in the continued fraction expansion of \(x\) are
bounded. Quadratic irrationals, including the golden ratio
\(\varphi=(1+\sqrt{5})/2\), belong to this class. Heuristically, such
numbers \textbf{avoid} exceptionally accurate low-height rationals: at
fixed \(Q\) they often sit closer to the \textbf{bulk} of
\(\delta_Q(X)\) than do numbers that admit a stellar approximant with
\(q\leq Q\) (e.g.~\(\pi\) near \(355/113\) when \(Q\ge 113\)).
\end{frame}

\begin{frame}{Transcendence versus fixed-\(Q\) floating-point
measurement}
\protect\phantomsection\label{transcendence-versus-fixed-q-floating-point-measurement}
Being transcendental excludes exact roots of nontrivial integer
polynomials, but it does \textbf{not} determine \(\delta_Q(x)\) at
moderate \(Q\) in double precision. Empirical comparison against
synthetic baselines therefore remains a useful \textbf{expository}
device: it shows how a single outstanding rational hit shifts
percentiles even when classical transcendentality proofs are orthogonal
to \eqref{eq:delta}.
\end{frame}

\begin{frame}[fragile]{Relation to the pooled reference}
\protect\phantomsection\label{relation-to-the-pooled-reference}
The Monte Carlo pipeline does not sample from Haar measure on
\(\mathbb{R}/\mathbb{Z}\) in a measure-theoretic sense; it mixes simple
laws on \([0,1)\) chosen for reproducibility and interpretability.
Percentiles are \textbf{always conditional} on that mixture. Changing
\texttt{experiment.n\_uniform}, \texttt{n\_quadratic}, \texttt{n\_beta},
or \texttt{quadratic\_candidates} shifts the empirical CDF;
\texttt{compare\_mod1:\ true} aligns constants with \([0,1)\) when
comparing primarily to the uniform component.
\end{frame}

\end{document}
